% ==================================================================
% id:       func-016
% topic:    פונקציות
% subtopic: פונקציה הפוכה
% solution: full
% ==================================================================
הראו כי הפונקציה
$\ds{f(x)=\frac{x+1}{x-1}}$
חד חד ערכית בתחום הגדרתה.

\hint{הניחו ש-$f(x_1)=f(x_2)$ והגיעו ל-$x_1=x_2$ (למשל בכפל בהצלבה).}
\begin{solbox}
תחום ההגדרה הוא $\mathbb{R}\setminus\{1\}$. יהיו $x_1,x_2$ בתחום ההגדרה ונניח כי $f(x_1)=f(x_2)$. אז
\[
.\frac{x_1+1}{x_1-1}=\frac{x_2+1}{x_2-1}
\]
כפל בהצלבה נותן
\[
,(x_1+1)(x_2-1)=(x_2+1)(x_1-1)
\]
ומכאן לאחר פתיחת סוגריים מתקבל $x_1=x_2$. לכן $f$ חד-חד ערכית.
\end{solbox}
